\section*{4.}
\subsection*{a.}
Es posible que dos departamentos tengan el nombre \textit{`Sistemas'} pues, como no son parte de la llave primaria (PK) del tipo de entidad, entonces no tiene restricciones de unicidad explícitas, por lo que la afirmación es verdadera.

\subsection*{b.}
Dada la estructura de la tabla de relación de \texttt{Trabajar}, cuya PK se compone de 1 entidad de \texttt{Empleado} y 1 de \texttt{Departamento}, entonces la restricción solamente indica que no se puede repetir una relación de trabajo entre los mismos valores de identificador de \texttt{Empleado}, de identificador de \texttt{Departamento} y del atributo \texttt{desde}, por lo que, si dos entidades de tipo \texttt{Empleado} son distintas, entonces sí se pueden relacionar con la misma entidad \texttt{Departamento}, por lo que la afirmación es verdadera.

\subsection*{c.}
Para las tablas de relación de \texttt{Trabajar} o de \texttt{Administrar}, el que haya identificadores de \texttt{Empleado} o de \texttt{Departamento} registrados en una fila de una de las tablas no genera restricciones en la otra a nivel lógico, ya que éstas no tienen relación entre sí, por lo que, como ambas tablas \texttt{Trabajar} y \texttt{Administrar} pueden tener filas con los mismos identificadores, entonces un empleado puede trabajar en un departamento y administrar otro a la vez, por lo que la afirmación es verdadera.

\subsection*{d.}
Como se explicó en el inciso anterior, \texttt{Trabajar} o de \texttt{Administrar} no tienen restricciones entre sí en cuanto a los valores que se pueden o deben almacenar por lo que, como un \texttt{Empleado} puede trabajar en un \texttt{Departamento} y administrar otro a la vez, entonces no requiere trabajar en el \texttt{Departamento} que administra para cumplir con las restricciones, por lo que la afirmación es falsa.

\subsection*{e.}
Ya que la PK de \texttt{Trabajar} determina que la tupla que la compone debe ser única, entonces en esta relación no se pueden repetir a la vez los tres valores \texttt{NúmEmpleado}, \texttt{NúmDepto} \texttt{Desde}. No obstante, esto no se restringe a la repetición de hasta dos valores (pues la PK sigue siendo identificable a partir de los diferentes valores del registro de lque no se repite), entonces se pueden repetir los valores de \texttt{NúmEmpleado} y \texttt{Desde} en una fila si se cambian los valores de \texttt{NúmDepto} i.e. un empleado puede trabajar en dos departamentos diferentes a partir de la misma fecha, por lo que la afirmación es verdadera.

\subsection*{f.}
Ya que puede un modelo relacional por si mismo no puede aplicar restricciones entre los atributos de una misma entidad, y que el diagrama no presenta espacificaciones adjuntas, entonces no se tiene una restricción entre \texttt{hasta} y \texttt{desde} que determine que uno debe ser anterior al otro, por lo que la afirmación es falsa.

\subsection*{g.}
Como se mencionó en el inciso $e.$, la relación \texttt{Trabajar} permite repetir valores de la PK siempre y cuando no sean todos. Es por esto que la tabla \texttt{Trabajar} refleja que la relación entre \texttt{NúmEmpleado} y \texttt{NúmDepto} es de cardinalidad $N:M$. Por esto, dado un identificador de empleado $x$, si bien es posible conocer los departamentos en los que trabaja el respectivo empleado mediante la línea de comandos
\begin{center}
    \texttt{SELECT NúmDepartamento FROM Trabajar WHERE NúmEmpleado = x;}
\end{center}
querer conocer exactamente el \textbf{único} \texttt{Departamento} donde éste trabaja no tiene sentido, por lo que la afirmación es falsa.

\subsection*{h.}
Si bien se pueden gestionar correctamente las restricciones entre valores en la implementación del diagrama, como se comentó en el inciso $f.$, el diagrama por sí mismo carece de dichas restricciones, por lo que, ya que la PK de la tabla de \texttt{Salario} se conforma por la dupla de \texttt{NúmEmpleado} y \texttt{Desde}, basta con que el segundo atributo sea distinto entre dos filas para que el empleado se pueda repetir por lo que, si se dan los siguientes condiciones en $Salario$, entonces un mismo empleado sí puede tener asignado más de un salario, por lo que la afirmación del inciso es falsa:
\begin{enumerate}
    \item La fecha de inicio no se repite pero el identificador de empleado sí en dos filas (asignación de salario al mismo empleado).
    \item El salario es distinto entre ambas filas (asignación de salarios distintos al empleado).
    \item Las fechas de \texttt{hasta} son posteriores a sus respectivas fechas de \texttt{desde} (se da intervalo real), y al menos un día se encuentra en ambos intervalos a la vez (los dos periodos de salario se dan a la vez).
\end{enumerate}

\subsection*{i.}
Ya que, por definición, todas las columnas que forman parte de la PK son estrictamente no nulas entonces, por cada registro de \texttt{Salario}, debe haber un identificador de empleado y una fecha de inicio asociados i.e. obligatoriamente debe haber valores para \texttt{NúmEmpleado} y \texttt{Desde} asociados, por lo que la proposición del inciso es es falsa.

\subsection*{j.}
Dada la cardinalidad expresada entre \texttt{Administrar} y \texttt{Departamento}, que indica una relación parcial para este último, queda explícito que un \texttt{Departamento} puede ser registrado en la base de datos sin necesidad de tener a alguien encargado de su administración, por lo que la afirmación es falsa.